\begin{table}[htbp]
\centering
\caption{Calibration quality comparison across weighting schemes}
\label{tab:calibration}
\renewcommand{\arraystretch}{1.2}
\resizebox{0.98\textwidth}{!}{%
\begin{tabular}{lrrrrr}
\toprule
& & \multicolumn{2}{c}{Equal Weighting ($W=I$)}
& \multicolumn{2}{c}{Vega Weighting} \\
\cmidrule(lr){3-4} \cmidrule(lr){5-6}
Underlying & $N$ &Stab. margin pass rate & RMSE (median)
& Stab. margin pass Rate & RMSE (median) \\
\midrule
SSE 50
& 13,373
& 78.81%
& $1.50\times10^{-4}$
& 93.56%
& $1.38\times10^{-4}$ \\

CSI 300
& 15,562
& 86.07%
& $1.09\times10^{-4}$
& 97.30%
& $1.02\times10^{-4}$ \\

CSI 1000
& 15,536
& 87.08%
& $1.37\times10^{-4}$
& 94.75%
& $1.34\times10^{-4}$ \\

\textbf{Overall}
& \textbf{44,471}
& \textbf{84.24%}
& $1.31\times10^{-4}$
& \textbf{95.28%}
& $1.23\times10^{-4}$ \\
\bottomrule
\end{tabular}
}
\begin{minipage}{0.98\textwidth}
\footnotesize
\vspace{0.6em}
\setstretch{1.1}
\footnotetext{
\textit{Note:} Stab. margin pass rate refers to the proportion of normal fitted slices with a nonpositive stability margin, i.e., slices whose fitted surface does not exhibit the branch-intersection distortion described in this section.}
\vspace{0.4em}
\end{minipage}
\end{table}
